\documentclass[12pt]{article}
\usepackage[T1]{fontenc}
\usepackage{lmodern}
\usepackage{amsmath,amssymb,amsthm,mathtools}
\usepackage{booktabs,array,longtable}
\usepackage[margin=1.05in]{geometry}
\usepackage{microtype}
\usepackage[hidelinks]{hyperref}
\usepackage{enumitem}
\usepackage{fancyhdr}
\usepackage{mathrsfs}
\allowdisplaybreaks[1]
\setcounter{tocdepth}{2}
\newtheorem{theorem}{Theorem}[section]
\newtheorem{lemma}[theorem]{Lemma}
\newtheorem{proposition}[theorem]{Proposition}
\newtheorem{corollary}[theorem]{Corollary}
\newtheorem{definition}[theorem]{Definition}
\theoremstyle{remark}
\newtheorem{remark}[theorem]{Remark}
\newtheorem{example}[theorem]{Example}
\newtheorem{exercise}[theorem]{Exercise}
\newcommand{\C}{\mathbb C}
\newcommand{\A}{\mathbb A}
\newcommand{\Hom}{\operatorname{Hom}}
\newcommand{\Ind}{\operatorname{Ind}}
\newcommand{\Res}{\operatorname{Res}}
\newcommand{\tr}{\operatorname{tr}}
\newcommand{\id}{\mathbf 1}
\newcommand{\cl}{\operatorname{Cl}}
\newcommand{\Irr}{\operatorname{Irr}}
\newcommand{\cG}{\mathcal C(G)}
\newcommand{\cH}{\mathcal C}
\newcommand{\Orb}{\mathcal O}
\newcommand{\eps}{\varepsilon}
\newcommand{\diag}{\operatorname{diag}}
\newcommand{\conv}{\mathbin{\widehat\otimes}}
\newcommand{\Gal}{\Gamma}
\pagestyle{fancy}
\fancyhf{}
\fancyhead[L]{\small Finite-group trace formulas}
\fancyhead[R]{\small Corrected and expanded notes}
\fancyfoot[C]{\thepage}
\setlength{\headheight}{14pt}
\title{\textbf{Finite-Group Trace Formulas:}\\
Compact Global and Local Analogies, and a $p$-Adic Bridge\\[5pt]
\large Corrected and expanded version of the 1996--2003 notes}
\author{David Joyner\\[4pt]
\normalsize Revised and expanded with mathematical and editorial assistance from ChatGPT (OpenAI)}
\date{Original notes: 1996--2003 \quad $\cdot$ \quad Revised: November 2026}
\begin{document}
\maketitle
\begin{abstract}
For a finite group $G$, the trace of a convolution operator on the
permutation representation $\C[\Gamma\backslash G]$ has two descriptions:
one involving conjugacy classes and orbital sums, the other involving
irreducible representations and their characters.  This elementary identity
models the Selberg trace formula for a compact quotient.  The analogous
bilinear identity on $\C[G]$ is a finite counterpart of the \emph{compact}
case of Arthur's local trace formula.  An important observation is that the
latter identity follows from the former by applying it to
$(G\times G,\Delta G)$, where $\Delta G$ is the diagonal subgroup.
We carefully fix counting-measure conventions, give proofs and explicit
examples for $S_3$ and $A_5$, and explain both the value and the limits
of the comparison with Arthur's global and local theories.
An additional section explains, through an exact $\mathrm{GL}_2(F)$
divergence calculation, why the noncompact $p$-adic local formula
requires truncation, Levi subgroups, and weighted distributions.
Throughout, we indicate how the finite-group formulations relate to
Arthur's finite example, Vign'eras's regular-bimodule method,
Yang's coset-space formulas, and Terras's finite spectral analogues.
A literature guide, glossary, exercises, and an account of corrections
appear at the end.
\end{abstract}
\noindent\textit{Historical and editorial note.}
The underlying exposition began as notes written by David Joyner during a
Stanford visit in summer 1996, with additions in 1997, 2000, and 2003.
The present revision corrects several representation-theoretic and
normalization errors, adds the diagonal-subgroup interpretation of the
local formula, and supplies background and references. The present edition also
adds a worked $p$-adic bridge. Its finite-group ideas have important
precedents in Arthur \cite{ArthurHecke1989}, Vign'eras
\cite{Vigneras1991}, and Terras \cite{Terras1999}; Yang's later
work \cite{Yang2006} treats closely related coset-space representations.
The original 1996--2003 notes predate Yang's 2006 publication. No
claim of priority is intended for the standard finite-group identities.
This is a revised exposition, not a reproduction of the original source.
\tableofcontents

\section{Why trace formulas have two sides}
\subsection{The global picture}
Let $H$ be a locally compact unimodular group and $\Gamma\subset H$ a
discrete subgroup such that $\Gamma\backslash H$ is compact.  Under
standard hypotheses on the test function $f$, right convolution on
$L^2(\Gamma\backslash H)$ has a kernel built from the translates of $f$.
Computing its trace by integrating the diagonal of the kernel yields
\emph{geometric terms}, grouped by conjugacy classes in $\Gamma$.
Decomposing the same operator into irreducible representations yields
\emph{spectral terms}.  In its schematic compact-quotient form,
\begin{equation}\label{eq:realcompact}
 \sum_{\pi}m_\Gamma(\pi)\,\tr\pi(f)
 =\sum_{[\gamma]_\Gamma}
 \operatorname{vol}(\Gamma_\gamma\backslash H_\gamma)
 \int_{H_\gamma\backslash H}f(x^{-1}\gamma x)\,d\dot x.
\end{equation}
Here $H_\gamma$ and $\Gamma_\gamma$ are centralizers, and the
measures must be chosen compatibly. Arthur's 1989 introductory
lecture derives this compact-quotient formula and then considers the
case of finite groups, obtaining Frobenius reciprocity as a special
case \cite[Example~2, p.~14]{ArthurHecke1989}; see also
\cite[\S1]{ArthurIntro2005}.
A finite group $G$, with any subgroup $\Gamma$, gives a setting in which
every sum is finite and none of the analytic convergence difficulties occur.

\subsection{The local picture}
For a local field $F$ and a connected reductive group $\mathbf G$ over
$F$, Arthur starts with the two-sided action of
$\mathbf G(F)\times\mathbf G(F)$ on $L^2(\mathbf G(F))$:
\begin{equation}\label{eq:two-sided-intro}
 ((a,b)\cdot\phi)(x)=\phi(a^{-1}xb).
\end{equation}
The trace of an operator obtained from two test functions can again be
analyzed geometrically and spectrally.  When $\mathbf G(F)$ is noncompact,
truncation is essential; the general result involves weighted orbital
integrals, weighted characters, Levi subgroups, and continuous spectral
parameters \cite{ArthurTowards1989,ArthurLocal1991}.
For a finite $G$, however, \eqref{eq:two-sided-intro} produces an
identity whose two sides involve nothing beyond orbital sums and
irreducible characters.  This is a model of the \emph{compact-group
case} of the local trace formula, not a formula for general reductive
$p$-adic groups. Vign'eras \cite[\S1.1]{Vigneras1991} uses exactly
this two-sided regular representation as the elementary entry point
to Arthur's theory, and develops the connection much further.

\subsection{Conventions and a map of the results}
\textbf{Every sum over $G$ in this note uses unnormalized counting
measure.} Thus $\sum_{g\in G}1=|G|$.  On $\Gamma\backslash G$ we also
use counting measure, so a delta function of one coset has norm one.
We shall make no tacit substitution between the normalized average
$|G|^{-1}\sum_G$ and the unnormalized sum $\sum_G$.
The main identities are:
\begin{align}
 \tr R_\Gamma(f)
 &=\sum_{\pi\in\Irr(G)}m_\Gamma(\pi)\,\tr\pi(f)
  =\sum_{[\gamma]_\Gamma}
    \frac{\Orb_G(f,\gamma)}{|\Gamma_\gamma|},
 \label{eq:preview-global}\\
 \sum_{\pi\in\Irr(G)}
   \tr\bar\pi(f_1)\,\tr\pi(f_2)
 &=\sum_{[\gamma]_G}
    \frac{\Orb_G(f_1,\gamma)\Orb_G(f_2,\gamma)}{|G_\gamma|}.
 \label{eq:preview-local}
\end{align}
The orbital sum $\Orb_G$ is defined below. The centralizers in the
two formulas are \emph{different}: $\Gamma_\gamma$ in the first,
$G_\gamma$ in the second.

\paragraph{Four complementary precedents.}
Arthur's finite-group example \cite[Example~2]{ArthurHecke1989}
explains the first formula through Frobenius reciprocity.
Vign\'eras \cite[\S1]{Vigneras1991} starts from the two-sided
$G\times G$ action underlying the second formula.
Yang \cite{Yang2006} develops related trace formulas for
coset-space representations, including vector-valued functions,
whereas Terras \cite[Chs.~22--23]{Terras1999} explores finite
Selberg-type identities in geometric and spectral settings.
These are overlapping but not identical approaches; the purpose
of the present note is to connect them using one fixed system
of normalization conventions.

\section{Finite harmonic analysis: notation and basic facts}
\subsection{Conjugacy and characters}
Write $\Irr(G)$ for the set of isomorphism classes of irreducible
finite-dimensional complex representations.  We denote the conjugacy
class of $\gamma\in G$ by
\[
 [\gamma]_G=\{x^{-1}\gamma x:x\in G\},\qquad
 G_\gamma=C_G(\gamma)=\{x\in G:x\gamma=\gamma x\}.
\]
The class has size $|G|/|G_\gamma|$.
We distinguish $[\gamma]_G$ from $[\gamma]_\Gamma$ when $\gamma\in\Gamma$.
They need not be the same set. Let $\cG=\{f:G\to\C\}$ and let
$\cG^G$ be the subspace of class functions, those constant on each
$G$-conjugacy class.
For $\pi\in\Irr(G)$, write $\chi_\pi(g)=\tr(\pi(g))$, and set
\[
 \pi(f)=\sum_{g\in G}f(g)\pi(g),\qquad
 \chi_\pi(f)=\tr\pi(f)=\sum_{g\in G}f(g)\chi_\pi(g).
\]
The contragredient (equivalently, complex-conjugate unitary)
representation is $\bar\pi$, whose character satisfies
$\chi_{\bar\pi}(g)=\overline{\chi_\pi(g)}=\chi_\pi(g^{-1})$.

\begin{theorem}[Character orthogonality and completeness]\label{thm:chars}
The irreducible characters are an orthonormal basis of $\cG^G$ for
\[
 \langle f,h\rangle_G
   =\frac1{|G|}\sum_{g\in G}f(g)\overline{h(g)}.
\]
In particular,
$|\Irr(G)|=\#\{G\text{-conjugacy classes}\}$.
For a treatment that derives orthogonality and completeness
systematically from the two-sided trace identity, see
\cite[\S\S1.1--1.9]{Vigneras1991}; compare the classical
character-theoretic approach in \cite{Serre}.
\end{theorem}
\begin{remark}
The characters are \emph{not} a basis of all of $\cG$ unless $G$ is
abelian.  Also, the ordinary point masses $\delta_g$ have squared norm
$1/|G|$ for the normalized inner product above: they are orthogonal,
but not normalized.  The functions $\sqrt{|G|}\,\delta_g$ are an
orthonormal basis of $\cG$.
\end{remark}

\subsection{Counting-measure orbital sums}
\begin{definition}
For $f\in\cG$ and $\gamma\in G$, the \emph{orbital sum} is
\begin{equation}\label{eq:def-orbit}
 \Orb_G(f,\gamma)=\sum_{x\in G}f(x^{-1}\gamma x).
\end{equation}
It is a class function in $\gamma$.
\end{definition}
Since each conjugate is attained by $|G_\gamma|$ elements $x$,
\begin{equation}\label{eq:orbitquot}
 \Orb_G(f,\gamma)
  =|G_\gamma|\sum_{\eta\in[\gamma]_G}f(\eta).
\end{equation}
The sum $\sum_{\eta\in[\gamma]_G}f(\eta)$ is the orbital integral
with counting measure on $G_\gamma\backslash G$.
Consequently the (finite) Weyl integration formula is simply
\begin{equation}\label{eq:weyl}
 \sum_{g\in G}f(g)
 =\sum_{[\gamma]_G}\frac{\Orb_G(f,\gamma)}{|G_\gamma|}.
\end{equation}
This is just a partition of $G$ into conjugacy classes.
It is the elementary precursor of the Weyl-integration
rearrangement used in Arthur's compact formula
\cite[\S1]{ArthurHecke1989} and in finite geometric analogues
such as \cite[Ch.~22]{Terras1999}.

\begin{proposition}[Fourier expansion of an orbital sum]\label{prop:fourier}
For every $f\in\cG$ and $\gamma\in G$,
\begin{equation}\label{eq:fourier}
 \boxed{\displaystyle
 \Orb_G(f,\gamma)=
 \sum_{\pi\in\Irr(G)}\chi_{\bar\pi}(f)\chi_\pi(\gamma).}
\end{equation}
\end{proposition}
\begin{proof}
Expand the class function $\gamma\mapsto\Orb_G(f,\gamma)$
in the orthonormal character basis. Its coefficient at $\chi_\pi$ is
\begin{align*}
 \frac1{|G|}\sum_{\gamma\in G}
       \Orb_G(f,\gamma)\overline{\chi_\pi(\gamma)}
 &=\frac1{|G|}\sum_{x,\gamma\in G}
       f(x^{-1}\gamma x)\overline{\chi_\pi(\gamma)}\\
 &=\sum_{y\in G}f(y)\overline{\chi_\pi(y)}
 =\chi_{\bar\pi}(f).
\end{align*}
In the second equality we used the conjugation-invariance of
$\chi_\pi$ and changed variables $y=x^{-1}\gamma x$.
\end{proof}
\begin{corollary}[Column orthogonality]\label{cor:column}
For $\eta,\gamma\in G$,
\begin{equation}\label{eq:columnorth}
 \sum_{\pi\in\Irr(G)}
   \overline{\chi_\pi(\eta)}\,\chi_\pi(\gamma)
 =\begin{cases}
    |G_\gamma|,&[\eta]_G=[\gamma]_G,\\
    0,&\text{otherwise}.
   \end{cases}
\end{equation}
\end{corollary}
\begin{proof}
Apply \eqref{eq:fourier} with $f=\delta_\eta$. Its left-hand side
counts the $x\in G$ for which $x^{-1}\gamma x=\eta$.
\end{proof}
\begin{remark}[The characteristic function of a class is not a point mass]
If $f=\mathbf1_{[\eta]_G}$, then
$\chi_\pi(f)=|[\eta]_G|\chi_\pi(\eta)$ and
$\Orb_G(f,\gamma)=|G|$ whenever $\gamma$ is conjugate to $\eta$.
If instead $f=\delta_\eta$, then $\chi_\pi(f)=\chi_\pi(\eta)$ and
the orbital sum in the matching class equals $|G_\eta|$.
Confusing these two functions loses a factor of the class size.
\end{remark}

\section{Two character-table examples}
\subsection{$S_3$}
The conjugacy classes of $S_3$ are represented by
$e$, $s=(12)$, and $t=(123)$; their sizes are $1,3,2$, and the
centralizer orders are $6,2,3$.  The irreducible representations are
the trivial $\id$, the sign representation $\eps$, and the
\emph{two-dimensional} standard representation $V$:
\[
 V=\{(z_1,z_2,z_3)\in\C^3:z_1+z_2+z_3=0\}.
\]
Equivalently, $V\cong\C^3/\C(1,1,1)$, but it is \emph{not}
$\C^3/V$, which is one-dimensional.
\begin{center}
\begin{tabular}{lrrr}
\toprule
 & $[e]$ (size $1$) & $[s]$ (size $3$) & $[t]$ (size $2$)\\
\midrule
$\chi_{\id}$ & $1$&$1$&$1$\\
$\chi_{\eps}$ & $1$&$-1$&$1$\\
$\chi_V$ & $2$&$0$&$-1$\\
\bottomrule
\end{tabular}
\end{center}
For example, $1^2+1^2+2^2=6=|G_e|$, while
$1^2+(-1)^2+0^2=2=|G_s|$, in agreement with
\eqref{eq:columnorth}.

\subsection{$A_5$}
In $A_5$ there are five conjugacy classes, with representatives
\[
 e,\quad a=(12)(34),\quad b=(123),\quad
 c=(12345),\quad d=c^2=(13524),
\]
and respective sizes $1,15,20,12,12$. Put
\[
 \phi_+=\frac{1+\sqrt5}{2},\qquad
 \phi_-=\frac{1-\sqrt5}{2}.
\]
The two numbers $\phi_+$ and $\phi_-$ are real algebraic
conjugates, \emph{not} complex conjugates.  With this notation:
\begin{center}
\begin{tabular}{lrrrrr}
\toprule
 & $e$ & $a$ & $b$ & $c$ & $d$\\
\midrule
class size & $1$&$15$&$20$&$12$&$12$\\
$\chi_1$ (degree 1) & $1$&$1$&$1$&$1$&$1$\\
$\chi_2$ (degree 3) & $3$&$-1$&$0$&$\phi_+$&$\phi_-$\\
$\chi_3$ (degree 3) & $3$&$-1$&$0$&$\phi_-$&$\phi_+$\\
$\chi_4$ (degree 4) & $4$&$0$&$1$&$-1$&$-1$\\
$\chi_5$ (degree 5) & $5$&$1$&$-1$&$0$&$0$\\
\bottomrule
\end{tabular}
\end{center}
As a useful check, the sum of the squares of the degrees is
$1+9+9+16+25=60=|A_5|$.  We will revisit the $3$-cycle class
when proving the local formula.

\section{The compact global trace formula for finite groups}
\subsection{The right-translation representation and its kernel}
Fix a subgroup $\Gamma\leq G$.  This is the finite
compact-quotient construction considered by Arthur
\cite[Example~2]{ArthurHecke1989}. For a closely related and
more general treatment using vector-valued functions on the coset
space, compare Yang \cite{Yang2006}.  Let
\[
 \cH_\Gamma=\C[\Gamma\backslash G]
 =\{\varphi:G\to\C:\varphi(\gamma x)=\varphi(x)
                    \text{ for }\gamma\in\Gamma\}.
\]
We use the honest right-translation \emph{representation}
\begin{equation}\label{eq:Rh}
 (R(g)\varphi)(x)=\varphi(xg),\qquad
 R(g_1)R(g_2)=R(g_1g_2).
\end{equation}
(Using $\varphi(xg^{-1})$ instead, with operators composed in the
usual order, gives an anti-representation.) For $f\in\cG$ define
\[
 R_\Gamma(f)=\sum_{g\in G}f(g)R(g).
\]
Choose counting measure on the coset set.  Then
\begin{equation}\label{eq:kernel}
 (R_\Gamma(f)\varphi)(x)
  =\sum_{y\in\Gamma\backslash G} K_f(x,y)\varphi(y),\quad
 K_f(x,y)=\sum_{\gamma\in\Gamma}f(x^{-1}\gamma y).
\end{equation}
The kernel is well-defined on left cosets in each variable:
replacing $x$ or $y$ by a representative of the same coset
simply reindexes the sum. In particular,
\begin{align}
 \tr R_\Gamma(f)
 &=\sum_{x\in\Gamma\backslash G} K_f(x,x)
  =\frac1{|\Gamma|}\sum_{x\in G}\sum_{\gamma\in\Gamma}
                                  f(x^{-1}\gamma x)\notag\\
 &=\frac1{|\Gamma|}\sum_{\gamma\in\Gamma}
                                  \Orb_G(f,\gamma).
 \label{eq:rawgeo}
\end{align}
The factor $1/|\Gamma|$ is the passage from a sum over elements
of $G$ to a sum over left cosets. It is \emph{not} a factor $1/|G|$.

\subsection{Geometric expansion: two valid ways to group terms}
Within $\Gamma$, group the summands of \eqref{eq:rawgeo} by
$\Gamma$-conjugacy classes.  The orbital sum
$\Orb_G(f,\gamma)$ is constant along each such class, and
$|[\gamma]_\Gamma|=|\Gamma|/|\Gamma_\gamma|$.
Hence
\begin{equation}\label{eq:geoGamma}
 \boxed{\displaystyle
 \tr R_\Gamma(f)=
 \sum_{[\gamma]_\Gamma}\frac{\Orb_G(f,\gamma)}{|\Gamma_\gamma|}.}
\end{equation}
Alternatively, one may group \emph{by $G$-conjugacy classes}, but
then one must count only the elements of each class that lie in
$\Gamma$:
\begin{equation}\label{eq:geoG}
 \tr R_\Gamma(f)=
 \frac1{|\Gamma|}\sum_{[\gamma]_G:\,[\gamma]_G\cap\Gamma\ne\varnothing}
 \bigl|[\gamma]_G\cap\Gamma\bigr|\Orb_G(f,\gamma).
\end{equation}
It is generally \emph{incorrect} to replace
$\bigl|[\gamma]_G\cap\Gamma\bigr|$ by
$\bigl|[\gamma]_G\bigr|$.

\begin{remark}[Comparison with the compact-quotient coefficient]
Define the quotient-orbit sum
\[
 \Orb_G^{\mathrm{quot}}(f,\gamma)
  =\sum_{x\in G_\gamma\backslash G}f(x^{-1}\gamma x)
  =\frac{\Orb_G(f,\gamma)}{|G_\gamma|}.
\]
Then \eqref{eq:geoGamma} becomes
\[
 \sum_{[\gamma]_\Gamma}
  \frac{|G_\gamma|}{|\Gamma_\gamma|}
  \Orb_G^{\mathrm{quot}}(f,\gamma).
\]
This is exactly the finite counting-measure counterpart of the
coefficient $\operatorname{vol}(\Gamma_\gamma\backslash H_\gamma)$
in \eqref{eq:realcompact}.
\end{remark}

\subsection{Spectral expansion and Frobenius reciprocity}
Over $\C$ the finite-dimensional representation $\cH_\Gamma$
is completely reducible:
\begin{equation}\label{eq:decomp}
 \cH_\Gamma\cong
 \bigoplus_{\pi\in\Irr(G)}m_\Gamma(\pi)\,\pi.
\end{equation}
It is the permutation module $\Ind_\Gamma^G\id_\Gamma$, so
Frobenius reciprocity gives
\begin{equation}\label{eq:mult}
 \boxed{\displaystyle
 m_\Gamma(\pi)=\dim\pi^\Gamma
  =\frac1{|\Gamma|}\sum_{\gamma\in\Gamma}\chi_\pi(\gamma).}
\end{equation}
Thus the multiplicity formula is precisely the finite-group
Frobenius-reciprocity calculation singled out by Arthur
\cite[p.~14]{ArthurHecke1989}. The averaging operator
$P_\Gamma=|\Gamma|^{-1}\sum_{\gamma\in\Gamma}\pi(\gamma)$
is the projection onto the $\Gamma$-fixed vectors. Taking its trace
proves the last equality directly. It follows that
\begin{equation}\label{eq:spec}
 \tr R_\Gamma(f)
 =\sum_{\pi\in\Irr(G)}m_\Gamma(\pi)\chi_\pi(f).
\end{equation}
The multiplicity $m_\Gamma(\pi)$ is a nonnegative integer;
no unexplained dimension or group-order factor belongs in
\eqref{eq:spec}.

\begin{theorem}[Finite compact-global trace formula]\label{thm:global}
For any finite group $G$, subgroup $\Gamma\leq G$, and function
$f:G\to\C$,
\begin{equation}\label{eq:globalboxed}
 \boxed{\displaystyle
 \sum_{\pi\in\Irr(G)}m_\Gamma(\pi)\chi_\pi(f)
 =\sum_{[\gamma]_\Gamma}
      \frac{\Orb_G(f,\gamma)}{|\Gamma_\gamma|}
 =\frac1{|\Gamma|}\sum_{\gamma\in\Gamma}
      \Orb_G(f,\gamma).}
\end{equation}
An equivalent Fourier expression is
\begin{equation}\label{eq:globalfourier}
 \tr R_\Gamma(f)
 =\frac1{|\Gamma|}\sum_{\pi\in\Irr(G)}
     \chi_{\bar\pi}(f)\sum_{\gamma\in\Gamma}\chi_\pi(\gamma).
\end{equation}
\end{theorem}
\begin{proof}
The first equality in \eqref{eq:globalboxed} equates
\eqref{eq:spec} and \eqref{eq:geoGamma}; the second is
\eqref{eq:rawgeo}. Inserting \eqref{eq:fourier} into
\eqref{eq:rawgeo} gives \eqref{eq:globalfourier}. The latter agrees
with the first expression because $m_\Gamma(\bar\pi)
=m_\Gamma(\pi)$ and $\pi\mapsto\bar\pi$ permutes $\Irr(G)$.
\end{proof}
\begin{example}[A diagnostic normalization check]\label{ex:s3global}
Let $G=S_3$, $\Gamma=\{e,(12)\}$, and $f=\delta_e$.
Then $R_\Gamma(e)$ acts as the identity on the three-dimensional
space $\C[\Gamma\backslash G]$, so its trace is $3$.
The geometric formula gives
\[
 \frac{\Orb_G(\delta_e,e)}{|\Gamma_e|}
 =\frac{6}{2}=3.
\]
And $\Ind_\Gamma^{S_3}\id\cong\id\oplus V$, hence the
spectral side is $1+2=3$. An expression with the coefficient
$|G_e|^{-1}$ in place of $|\Gamma_e|^{-1}$ would incorrectly give $1$.
\end{example}

\subsection{Special cases and correct interpretations}
\begin{corollary}[The regular representation: $\Gamma=\{e\}$]
When $\Gamma=\{e\}$, $m_\Gamma(\pi)=\dim\pi$, and
\begin{equation}\label{eq:reg}
 \sum_{\pi\in\Irr(G)}(\dim\pi)\chi_\pi(f)
  =|G|\,f(e).
\end{equation}
Equivalently, $f(e)=|G|^{-1}\sum_\pi (\dim\pi)\chi_\pi(f)$.
This is the finite-group Fourier inversion formula at the identity.
\end{corollary}
\begin{corollary}[The one-dimensional quotient: $\Gamma=G$]
When $\Gamma=G$, the representation $\C[G\backslash G]$ is trivial:
\begin{equation}\label{eq:gammaG}
 \sum_{g\in G}f(g)
 =\sum_{[\gamma]_G}\frac{\Orb_G(f,\gamma)}{|G_\gamma|}.
\end{equation}
The right-hand side is the finite Weyl integration formula
\eqref{eq:weyl}. The spectral side has \emph{only the trivial
representation}; it is not the full finite-group Plancherel sum.
\end{corollary}
\begin{corollary}[An index-two subgroup]\label{cor:index2}
If $\Gamma\lhd G$ has index two, and
$\eps:G\to\{\pm1\}$ is the nontrivial character of $G/\Gamma$,
then
\[
 \C[\Gamma\backslash G]\cong\id\oplus\eps,
 \qquad
 \sum_{[\gamma]_\Gamma}
     \frac{\Orb_G(f,\gamma)}{|\Gamma_\gamma|}
 =\chi_{\id}(f)+\chi_\eps(f)
 =2\sum_{\gamma\in\Gamma}f(\gamma).
\]
\end{corollary}
\begin{proof}
The two cosets are permuted through the quotient $G/\Gamma\cong C_2$;
its permutation representation is $\id\oplus\eps$.
\end{proof}

\subsection{A caution about induced representations}
The character of $\Ind_\Gamma^G\sigma$, for a representation
$\sigma$ of $\Gamma$, is
\begin{equation}\label{eq:indchar}
 \chi_{\Ind_\Gamma^G\sigma}(g)
  =\frac1{|\Gamma|}
    \sum_{\substack{x\in G\\x^{-1}gx\in\Gamma}}
       \chi_\sigma(x^{-1}gx).
\end{equation}
The factor $1/|\Gamma|$ is essential.  In general the argument
``$x^{-1}\Gamma x=\Gamma$'' is false: it holds only if $\Gamma$
is normal (or for $x$ in its normalizer).  Thus one cannot conclude
that the induced representation has no $\Gamma$-fixed vectors just
because $\sigma\not\cong\id_\Gamma$.

For example, with $G=S_3$, $\Gamma=\langle(12)\rangle$, and
$\sigma=\eps_\Gamma$ the nontrivial character of $C_2$, one has
\[
 \Ind_\Gamma^G\sigma\cong\eps_G\oplus V.
\]
The restriction $V|_\Gamma\cong\id_\Gamma\oplus\eps_\Gamma$,
so $\Ind_\Gamma^G\sigma$ \emph{does} have a
$\Gamma$-fixed vector.  For a normal subgroup $\Gamma\lhd G$,
in contrast, the restriction of $\Ind_\Gamma^G\sigma$ is a sum
of conjugates of $\sigma$, and a nontrivial irreducible $\sigma$
then gives no $\Gamma$-fixed vectors.  The appropriate general tool
is Mackey's double-coset formula, not the unqualified vanishing claim;
see \cite[Chapters 7--8]{Serre}. Yang's coset-space trace formulas
\cite{Yang2006} provide a complementary framework for extending
the scalar permutation-module setting used here.

\section{The compact local trace formula for finite groups}
\subsection{The bilinear identity}
Let $f_1,f_2:G\to\C$, and consider the finite sum
\begin{equation}\label{eq:Idef}
 I(f_1,f_2)=\sum_{x,y\in G}
          f_1(x^{-1}yx)f_2(y).
\end{equation}
This is the two-sided trace computation appearing, in different
notation, in Vign\'eras \cite[\S1.1]{Vigneras1991}.
Using the orbital expansion \eqref{eq:fourier}, we obtain
\begin{align}
 I(f_1,f_2)
 &=\sum_{y\in G}\Orb_G(f_1,y)f_2(y)\notag\\
 &=\sum_{\pi\in\Irr(G)}
     \chi_{\bar\pi}(f_1)\chi_\pi(f_2).
 \label{eq:local-spec}
\end{align}
On the other hand, partition the $y$-sum by conjugacy classes.  For
any class function $A(y)$,
$\sum_y A(y)f_2(y)=
\sum_{[\gamma]_G}A(\gamma)\Orb_G(f_2,\gamma)/|G_\gamma|$.
Taking $A(y)=\Orb_G(f_1,y)$ gives
\begin{equation}\label{eq:local-geom}
 I(f_1,f_2)
 =\sum_{[\gamma]_G}\frac{\Orb_G(f_1,\gamma)
                               \Orb_G(f_2,\gamma)}{|G_\gamma|}.
\end{equation}
\begin{theorem}[Finite compact-local trace formula]\label{thm:local}
For arbitrary complex-valued functions $f_1,f_2$ on $G$,
\begin{equation}\label{eq:localboxed}
 \boxed{\displaystyle
 \sum_{\pi\in\Irr(G)}
  \chi_{\bar\pi}(f_1)\chi_\pi(f_2)
 =\sum_{[\gamma]_G}\frac{\Orb_G(f_1,\gamma)
                             \Orb_G(f_2,\gamma)}{|G_\gamma|}.}
\end{equation}
\end{theorem}
\begin{proof}
Equate \eqref{eq:local-spec} and \eqref{eq:local-geom},
which both compute \eqref{eq:Idef}.
\end{proof}
\begin{remark}[What Vign\'eras's treatment adds]
Vign\'eras \cite[\S1]{Vigneras1991} begins with a field
$L$ whose characteristic does not divide $|G|$, rather than
only $\C$, and derives Schur's lemma, completeness, and
orthogonality from this starting point. We retain $\C$
so that contragredients, complex conjugation, and the positive
identity \eqref{eq:squares} have their familiar Hilbert-space
meaning. In characteristic dividing $|G|$, semisimplicity fails
in general and the subsequent spectral argument requires a
different formulation.
\end{remark}
\begin{corollary}[Positive form; class-function Plancherel]\label{cor:squares}
For every complex-valued $f$ on $G$,
\begin{equation}\label{eq:squares}
 \boxed{\displaystyle
 \sum_{\pi\in\Irr(G)}|\chi_\pi(f)|^2
 =\sum_{[\gamma]_G}
      \frac{|\Orb_G(f,\gamma)|^2}{|G_\gamma|}.}
\end{equation}
\end{corollary}
\begin{proof}
In \eqref{eq:localboxed}, put $f_1=f$ and $f_2=\bar f$.
The left-hand side becomes
$\sum_\pi|\chi_{\bar\pi}(f)|^2$, which equals
$\sum_\pi|\chi_\pi(f)|^2$ after relabeling irreducibles;
$\Orb_G(\bar f,\gamma)=\overline{\Orb_G(f,\gamma)}$.
\end{proof}
\begin{remark}
This identity is a Parseval/Plancherel statement for the
\emph{conjugacy-invariant projection} of $f$, with the particular
unnormalized Fourier and orbital transforms chosen here. It does not
assert that character traces determine an arbitrary function on $G$:
they determine its conjugation average.
\end{remark}

\subsection{A stronger interpretation: the diagonal subgroup}
The preceding two identities are not unrelated analogies.  There is an
exact reduction of the \emph{finite local} identity to the \emph{finite
global} identity. The regular $G\times G$ representation itself
is the point of departure in Vign\'eras \cite[\S1.1]{Vigneras1991};
our additional point here is to obtain it by applying the
coset-space theorem to the diagonal subgroup.
\begin{theorem}[The local identity is a global identity on
$G\times G$]\label{thm:diagonal}
Let $H=G\times G$, let
$\Delta G=\{(g,g):g\in G\}\leq H$, and define
$F(a,b)=f_1(a)f_2(b)$.  Applying
Theorem~\ref{thm:global} to $(H,\Delta G,F)$ gives exactly
Theorem~\ref{thm:local}.
\end{theorem}
\begin{proof}
The map
\[
 \Delta G\backslash(G\times G)\longrightarrow G,
 \qquad [(g_1,g_2)]\longmapsto g_1^{-1}g_2
\]
is a bijection.  Under it, the right action of $(a,b)\in G\times G$
becomes
\begin{equation}\label{eq:two-sided}
 ((a,b)\cdot\varphi)(x)=\varphi(a^{-1}xb),
\end{equation}
Arthur's two-sided regular action.  Finite Peter--Weyl theory
(or decomposition of the regular bimodule) gives
\begin{equation}\label{eq:peter-weyl}
 \C[G]\cong
   \bigoplus_{\pi\in\Irr(G)}\bar\pi\boxtimes\pi
 \qquad\text{as a $(G\times G)$-module}.
\end{equation}
Each summand occurs once.  Its character at $(a,b)$ is
$\chi_{\bar\pi}(a)\chi_\pi(b)$. Thus the spectral side for $F$ is
$\sum_\pi\chi_{\bar\pi}(f_1)\chi_\pi(f_2)$.

Conjugacy classes of $\Delta G$ are indexed by $G$-conjugacy classes:
$[(\gamma,\gamma)]_{\Delta G}$ corresponds to $[\gamma]_G$.
Its centralizer in $\Delta G$ has order $|G_\gamma|$.  Moreover
\[
 \Orb_{G\times G}(F,(\gamma,\gamma))
 =\sum_{x,y\in G}f_1(x^{-1}\gamma x)f_2(y^{-1}\gamma y)
 =\Orb_G(f_1,\gamma)\Orb_G(f_2,\gamma).
\]
The global geometric expansion \eqref{eq:geoGamma}, now for
$(H,\Delta G)$, is exactly the right-hand side of
\eqref{eq:localboxed}.  This proves the assertion.
\end{proof}
\begin{remark}[A kernel matching Arthur's construction]
The convolution operator for $F=f_1\otimes f_2$ on $\C[G]$ is
\[
 (T_F\varphi)(x)=
 \sum_{a,b\in G}f_1(a)f_2(b)\varphi(a^{-1}xb).
\]
Its matrix kernel (with counting measure in the $y$ variable) is
\[
 K_F(x,y)=\sum_{a\in G}f_1(a)f_2(x^{-1}ay),
\]
so
\[
 \tr T_F=\sum_{x,a\in G}f_1(a)f_2(x^{-1}ax)=I(f_1,f_2).
\]
The last equality follows by interchanging the two functions and
changing variables in \eqref{eq:Idef}.  This double-sided kernel is a
closer parallel to Arthur's local trace-formula starting point than
the bilinear identity viewed solely as character orthogonality;
compare \cite[\S1.1]{Vigneras1991} and
\cite[\S1]{ArthurTowards1989}.
\end{remark}

\subsection{The character-table matrix is orthogonal after weighting}
Choose representatives $\gamma_1,\ldots,\gamma_r$ for
$G$-conjugacy classes and irreducibles
$\pi_1,\ldots,\pi_r$; set $c_j=|G_{\gamma_j}|$.  For each $j$ let
\begin{equation}\label{eq:basis-orb}
 f_j=|G|^{-1}\mathbf1_{[\gamma_j]_G}.
\end{equation}
Then $\Orb_G(f_j,\gamma_i)=\delta_{ij}$, by
\eqref{eq:def-orbit}. Define
\[
 a_{ij}=\chi_{\pi_i}(f_j)
 =\frac{\chi_{\pi_i}(\gamma_j)}{c_j},\qquad
 A=(a_{ij})_{1\leq i,j\leq r}.
\]
The local formula for $f_i,f_j$ gives column orthogonality:
\begin{equation}\label{eq:matrix}
 \boxed{\displaystyle
 A^*A=\diag(c_1^{-1},\ldots,c_r^{-1}).}
\end{equation}
Equivalently,
$A\diag(\sqrt{c_1},\ldots,\sqrt{c_r})$ is unitary.
\emph{The order of multiplication in \eqref{eq:matrix} matters.}
The matrix $AA^*$ is generally not diagonal.
\begin{example}[$S_3$ revisited]
With $\gamma_1=e$, $\gamma_2=(12)$, $\gamma_3=(123)$,
\[
 A=\begin{pmatrix}
 1/6&1/2&1/3\\
 1/6&-1/2&1/3\\
 1/3&0&-1/3
 \end{pmatrix}.
\]
A direct calculation yields
\[
 A^*A=\diag(1/6,1/2,1/3),\qquad
 (AA^*)_{12}=-\frac19\ne0.
\]
Thus a claim that $AA^*$ is diagonal would be false even for $S_3$.
\end{example}
\begin{example}[$A_5$ and a class indicator]\label{ex:A5}
Let $f=\mathbf1_{[(123)]_{A_5}}$, the indicator of the class of
$20$ three-cycles. From the character table,
\[
 (\chi_1(f),\ldots,\chi_5(f))=(20,0,0,20,-20),
\]
so the left side of \eqref{eq:squares} is
$20^2+20^2+(-20)^2=1200$.  The orbital sum is zero on other
classes and is $60$ at a three-cycle, since $|A_5|=60$.
Its centralizer has order $3$, so the right-hand side is
$60^2/3=1200$.  The equality holds with no hidden measure factors.
\end{example}

\section{What these finite identities do and do not model}
\subsection{The compact-quotient Selberg formula}
The finite global formula, Theorem~\ref{thm:global}, reproduces the
\emph{form} of the Selberg trace formula for compact quotient:
this is the analogy made explicit in Arthur's finite-group example
\cite[Example~2]{ArthurHecke1989}.
\begin{center}
\begin{tabular}{@{}p{0.32\linewidth}p{0.54\linewidth}@{}}
\toprule
Finite group & Locally compact compact quotient\\
\midrule
$\C[\Gamma\backslash G]$ & $L^2(\Gamma\backslash H)$\\
$\chi_\pi(f)$ & trace of $\pi(f)$\\
$m_\Gamma(\pi)$ & discrete spectral multiplicity\\
$\Orb_G(f,\gamma)/|G_\gamma|$ & orbital integral over $H_\gamma\backslash H$\\
$|G_\gamma|/|\Gamma_\gamma|$ & $\operatorname{vol}(\Gamma_\gamma\backslash H_\gamma)$\\
finite sum on $\Gamma\backslash G$ & integral of the kernel on the diagonal\\
\bottomrule
\end{tabular}
\end{center}
Terras \cite[Ch.~22]{Terras1999} takes a different, complementary
route from trace formulas to finite spectral geometry, including
isospectral nonisomorphic examples; Chapter~23 turns to finite
upper half planes. Those applications require extra geometric
structures and should not be identified with the bare coset-space
identity above.

The compact quotient condition should not be confused with
compactness of the group itself.  For example, a cocompact lattice
in a noncompact Lie group yields a compact quotient.  Finite groups
are a particularly simple case in which both the group and the
quotient are compact.

\subsection{Why Arthur's global formula is harder}
In adelic examples, $H=\mathbf G(\A_F)$ (with the appropriate
central quotient) and $\Gamma=\mathbf G(F)$ for a reductive group
over a number field $F$. The quotient usually is not compact.
The kernel need not be integrable along its diagonal, and the spectrum
need not be purely discrete.  Selberg's and Arthur's theory handles
continuous spectrum, parabolic subgroups, Eisenstein series, and
truncation; refined formulas feature weighted geometric and spectral
terms. None of these difficulties can be detected in the finite
counting-measure identity. See
\cite{Gelbart1996,ArthurIntro2005,Arthur1978,Arthur1980}.

\subsection{Why the local analogy is genuinely relevant}
Arthur explicitly begins his 1989 article \emph{Towards a local trace
formula} with the two-sided $\mathbf G(F)\times\mathbf G(F)$-action,
then works out the compact-group case by a Weyl integration formula
and Peter--Weyl decomposition \cite[\S\S1--2]{ArthurTowards1989}.
Theorem~\ref{thm:diagonal} gives the precise finite-group version
of that construction. Vign\'eras \cite[\S\S1--2]{Vigneras1991}
provides an especially useful intermediate presentation from the
finite calculation to the $p$-adic analytic issues.
Section~\ref{sec:padic-bridge} analyzes
the noncompact transition in detail. When $\mathbf G(F)$ is noncompact, the
integral of the diagonal kernel may diverge. Arthur's full local trace
formula replaces the simple finite sums by a carefully truncated
identity whose geometric and spectral expansions include terms
indexed by Levi subgroups \cite{ArthurLocal1991}.
In particular, the finite formula is \emph{not} the general
Arthur local trace formula with all analytic difficulties suppressed:
its lack of proper parabolic/Levi data and continuous spectrum means
that major structural features are absent.


\section{From finite groups to reductive $p$-adic groups:\\
why the local trace formula acquires new terms}
\label{sec:padic-bridge}
The finite identity \eqref{eq:localboxed} has two ingredients:
conjugation orbits on the geometric side and a decomposition of the
regular representation on the spectral side. Both ingredients survive
for a reductive $p$-adic group.  What fails is the final, seemingly
innocent step of integrating the kernel along the diagonal.  This
section explains that failure and the resulting changes in the
formulas. The goal is a conceptual bridge, \emph{not} a proof of
Arthur's local trace formula; the latter requires substantial analytic
and combinatorial arguments \cite{ArthurTowards1989,ArthurLocal1991}.
Vign\'eras \cite[\S2]{Vigneras1991} offers a shorter
complementary orientation, especially for elliptic elements,
weighted orbital integrals, and tempered distributions.

\subsection{The setting: Haar measure, Hecke algebra, and admissibility}
Let $F$ be a non-Archimedean local field of characteristic zero, with
ring of integers $\mathcal O_F$, uniformizer $\varpi$, and residue
field of cardinality $q$.  Thus $|\varpi|_F=q^{-1}$.  Let
$\mathbf G$ be a connected reductive algebraic group over $F$, and
write $G=\mathbf G(F)$.  The group $G$ is locally compact,
totally disconnected, and unimodular.  Its Haar measure $dg$ is an
analogue of counting measure, but $\operatorname{vol}(G)$ is usually
infinite.  We use
\[
 \mathcal H(G)=C_c^\infty(G)
 =\{f:G\longrightarrow\C:\ f\text{ is locally constant and
 compactly supported}\},
\]
where ``smooth'' in this context means locally constant. For a
compact open subgroup $K\subset G$, $\mathbf 1_K$ belongs to
$\mathcal H(G)$; we often normalize $dg$ by
$\operatorname{vol}(K)=1$.

The analogue of a finite-dimensional complex representation is now
often a \emph{smooth admissible} representation $(\pi,V)$:
for every $v\in V$ the stabilizer of $v$ is open, and for every
compact open $K$ the space $V^K$ is finite-dimensional. For such a
representation,
\[
 \pi(f)=\int_G f(g)\pi(g)\,dg
\]
is finite-rank when $f\in\mathcal H(G)$: choose $K$ fixing $f$ on
both sides; then $\pi(f)$ has image in $V^K$. Its trace
$\Theta_\pi(f)=\operatorname{tr}\pi(f)$ is the
\emph{distribution character}. Unlike a finite-group character,
$\Theta_\pi$ need not be a globally defined ordinary function,
although Harish-Chandra's regularity theory represents it by a
locally integrable function and by a locally constant function on
the regular semisimple locus.  For foundational material see
\cite{Casselman1995,Silberger1979}.

\begin{remark}[A change of meaning for the diagonal subgroup]
For finite $G$, the identification
$\Delta G\backslash(G\times G)\simeq G$ in
Theorem~\ref{thm:diagonal} is an instance of a quotient by a
\emph{discrete} subgroup, and the quotient is compact. For a
non-discrete $p$-adic $G$, the same identification remains valid
as an isomorphism of homogeneous spaces, but $\Delta G$ is
\emph{closed and non-discrete}, while the quotient is usually
noncompact. One therefore cannot apply the usual compact-lattice
Selberg trace formula to $(G\times G,\Delta G)$ and expect a
local trace formula for free. The regular-bimodule analogy persists;
the trace-class argument does not.
\end{remark}

\subsection{The same kernel, but potentially an infinite trace}
Let $f_1,f_2\in\mathcal H(G)$, and use the same two-sided
representation as in \eqref{eq:two-sided-intro}:
\[
 (R(a,b)\phi)(x)=\phi(a^{-1}xb), \qquad \phi\in L^2(G).
\]
For $f=f_1\otimes f_2$, the operator
$R(f)=\int_{G\times G}f_1(a)f_2(b)R(a,b)\,da\,db$
has the integral kernel
\begin{equation}\label{eq:padic-kernel}
 K_f(x,y)=\int_G f_1(xu)f_2(uy)\,du.
\end{equation}
Indeed, substitute $a=xu$ and then $b=uy$ in the expression for
$R(f)\phi(x)$.  On the diagonal, the change of variables
$u\mapsto x^{-1}u$ yields
\begin{equation}\label{eq:padic-diagonal}
 K_f(x,x)=\int_G f_1(u)f_2(x^{-1}ux)\,du.
\end{equation}
For a finite group, integrating \eqref{eq:padic-diagonal} over $x$
gives exactly $I(f_1,f_2)$ of \eqref{eq:Idef}. For a compact group,
Haar integration and Peter--Weyl theory similarly give a well-defined
trace. For noncompact $G$, the tempting expression
\[
 \int_G K_f(x,x)\,dx
\]
can diverge.  It may still diverge after the obvious noncompactness
of the split center has been removed.  The following calculation
shows this explicitly.

\subsection{An exact divergence calculation for $\mathrm{GL}_2(F)$}
\begin{proposition}[Every Cartan shell contributes one unit]
\label{prop:gl2-divergence}
Take $G=\mathrm{GL}_2(F)$,
$K=\mathrm{GL}_2(\mathcal O_F)$, and
$f_1=f_2=\mathbf1_K$, with $\operatorname{vol}(K)=1$.
Let $Z=F^\times I_2$ be the center and give $Z\backslash G$
Haar measure for which the image $\bar K$ of $K$ has volume one.
Then $K_f(x,x)$ descends to $Z\backslash G$, but
\begin{equation}\label{eq:gl2-divergence}
 \int_{Z\backslash G}K_f(x,x)\,d\bar x=+\infty.
\end{equation}
More precisely, with $a_n=\operatorname{diag}(\varpi^n,1)$,
$\bar a_n\in Z\backslash G$, and
\[
 s_0=1,\qquad s_n=(q+1)q^{n-1}\quad(n\geq1),
\]
the Cartan double coset $\bar K\bar a_n\bar K$ has volume $s_n$,
while $K_f(x,x)=s_n^{-1}$ everywhere on that coset. Hence the
integral over the first $N+1$ double cosets is exactly $N+1$.
\end{proposition}
\begin{proof}
Formula \eqref{eq:padic-diagonal} gives
\begin{equation}\label{eq:gl2-overlap}
 K_f(x,x)=\int_G \mathbf1_K(u)\mathbf1_K(x^{-1}ux)\,du
 =\operatorname{vol}(K\cap xKx^{-1}).
\end{equation}
This is unchanged when $x$ is multiplied by a scalar matrix or
on either side by an element of $K$.  The Cartan decomposition
\[
 Z\backslash G=\coprod_{n\geq0}\bar K\bar a_n\bar K
\]
therefore partitions the quotient into shells on which
\eqref{eq:gl2-overlap} is constant.  For $n\geq1$ the index
\[
 [K:K\cap a_nKa_n^{-1}]
  =\#\mathbb P^1(\mathcal O_F/\varpi^n\mathcal O_F)
  =(q+1)q^{n-1}
\]
is $s_n$; equivalently this is the number of vertices at distance
$n$ from a fixed vertex in the $(q+1)$-regular Bruhat--Tits tree of
$\mathrm{PGL}_2(F)$.  Thus the intersection has volume $s_n^{-1}$,
whereas the associated double coset in $Z\backslash G$ has
volume $s_n$.  Each shell contributes $s_n s_n^{-1}=1$.
Summing over all $n$ proves \eqref{eq:gl2-divergence}.
\end{proof}
\begin{remark}[What the example proves---and does not prove]
This is an actual divergence of the \emph{untruncated diagonal
integral}, even modulo the split center $Z=A_G(F)$. It is not a
counterexample to Arthur's local trace formula: the latter is built
from a truncated kernel and a carefully extracted distribution.
The elementary radial cutoff
$u_N=\mathbf1_{\coprod_{n=0}^N\bar K\bar a_n\bar K}$
yields $\int_{Z\backslash G}K_f(x,x)u_N(x)\,d\bar x=N+1$,
which illustrates how polynomial growth in a truncation parameter
can arise.  Arthur's precise truncation is more structured than this
radial cutoff; we do \emph{not} identify their constant terms.
\end{remark}

\subsection{Why the spectral side also changes}
Finite Peter--Weyl theory decomposes $\C[G]$ into a finite direct
sum of $\bar\pi\boxtimes\pi$. In contrast, the unitary regular
representation of a reductive $p$-adic group is described by
Harish-Chandra's Plancherel theorem, involving direct \emph{integrals}
and a continuous family of tempered representations. One version of
Fourier inversion has the schematic form
\begin{equation}\label{eq:padic-plancherel}
 f(e)=\int_{\widehat G_{\mathrm{temp}}}
    \Theta_\pi(f)\,d\mu_{\mathrm{Pl}}(\pi),
 \qquad f\in\mathcal H(G),
\end{equation}
with the usual interpretation of the measure, characters, and
normalizations.  Unlike \eqref{eq:reg}, this is generally not a
finite character sum.

One source of the continuous parameters is \emph{normalized parabolic
induction}. For $G=\mathrm{GL}_2(F)$, let $B=TN$ be the upper
triangular Borel subgroup, with
\[
 T=\left\{\begin{pmatrix}a&0\\0&d\end{pmatrix}:
 a,d\in F^\times\right\},\quad
 N=\left\{\begin{pmatrix}1&u\\0&1\end{pmatrix}:
 u\in F\right\}.
\]
The modulus character is
$\delta_B(\operatorname{diag}(a,d))=|a/d|_F$.
For quasicharacters $\chi_1,\chi_2$ of $F^\times$, the normalized
principal-series representation is
\begin{equation}\label{eq:principal-series}
 \operatorname{Ind}_{B}^{G}
  (\chi_1\boxtimes\chi_2)
 =\left\{\phi:G\to\C:\ 
 \phi(tng)=\delta_B(t)^{1/2}
       \chi_1(a)\chi_2(d)\phi(g),\quad
 t=\operatorname{diag}(a,d),\ n\in N\right\},
\end{equation}
with the right-translation action and the usual smoothness
conditions.  Taking $\chi_j(x)=|x|_F^{it_j}$ already gives
continuous unitary parameters; reducibility and equivalences must
be treated separately.  Intertwining operators relate different
parabolic realizations. Their normalized forms, and derivatives
thereof, enter the weighted spectral distributions of the local
trace formula \cite{Casselman1995,ArthurLocal1991}.

\subsection{Levi subgroups and the geometric indexing}
A parabolic subgroup $P$ has a Levi decomposition $P=MN$:
$M$ is reductive and $N$ its unipotent radical. Fix a minimal
$F$-Levi subgroup $M_0$ and denote by $\mathcal L(M_0)$ the
Levi subgroups containing it.  For each $M$ let $A_M$ be the
maximal $F$-split torus in its center. These tori provide the
real vector spaces $\mathfrak a_M$ used in truncation.

For $\mathrm{GL}_2$ there are only two standard Levi choices:
\[
 \mathcal L(M_0)=\{M_0=\mathrm{GL}_1\times\mathrm{GL}_1,
                     \ G=\mathrm{GL}_2\}.
\]
A regular semisimple element whose characteristic polynomial splits
over $F$ is conjugate into $M_0$; it is elliptic in $M_0$ but not
in $G$.  A regular element with irreducible quadratic characteristic
polynomial has a centralizer isomorphic to $E^\times$ for a
quadratic field extension $E/F$, and $E^\times/F^\times$ is compact:
it is \emph{elliptic modulo the center} in $G$.
This elementary dichotomy is the rank-one shadow of the
Levi-indexed expansion on Arthur's geometric side.

For a regular semisimple $\gamma$ the ordinary orbital integral is
\begin{equation}\label{eq:padic-orbital}
 O_\gamma(f)=\int_{G_\gamma\backslash G}
       f(x^{-1}\gamma x)\,d\dot x,
\end{equation}
with compatible Haar measures.  In many conventions a
\emph{normalized} orbital integral includes the factor
$|D^G(\gamma)|_F^{1/2}$, where
\[
 D^G(\gamma)=\det(1-\operatorname{Ad}(\gamma)
                     \mid\mathfrak g/\mathfrak g_\gamma).
\]
This analytic discriminant is not the finite-group centralizer
factor $|G_\gamma|$ from \eqref{eq:localboxed}.

\subsection{Truncation, weighted orbital integrals, and weighted characters}
Write $A_G$ for the maximal split central torus of $G$.  Arthur
introduces a truncation function $u(x,T)$ on $A_G(F)\backslash G$,
parametrized by a point $T$ sufficiently far into a positive chamber,
and first studies
\begin{equation}\label{eq:padic-trunc}
 K^T(f)=\int_{A_G(F)\backslash G}
         K_f(x,x)u(x,T)\,d\bar x.
\end{equation}
In the non-Archimedean case the parameter lies on an appropriate
lattice. The truncation is designed to be computable in both
the geometric expansion (from Weyl integration) and the spectral
expansion (from Plancherel theory).

For a proper Levi $M$, the geometric expansion involves a weight
$v_M$, built from the relative positions of the parabolic height
functions $H_P(x)$ for parabolics with Levi factor $M$.
These heights lie in $\mathfrak a_M$ and encode the different
ways $x$ approaches infinity along a parabolic direction.  An
\emph{illustrative}, single-variable weighted orbital integral has
shape
\begin{equation}\label{eq:weighted-model}
 \int_{G_\gamma\backslash G}
    f(x^{-1}\gamma x)\,v_M(x)\,d\dot x,
\end{equation}
possibly multiplied by a Weyl discriminant or other measure factors.
In Arthur's local formula the functions are on $G\times G$,
and the actual weight depends on the pair $(x_1,x_2)$; the final
weighted distribution is a \emph{double} orbital integral. For
example, in the conventions of \cite[\S12]{ArthurLocal1991}, for a
suitable $G$-regular $\gamma\in M(F)$ it has the form
\begin{equation}\label{eq:arthur-weighted-pair}
 J_M(\gamma,f)=|D^G(\gamma)|_F
 \int_{A_M(F)\backslash G}\!\int_{A_M(F)\backslash G}
 f_1(x_1^{-1}\gamma x_1)f_2(x_2^{-1}\gamma x_2)
 v_M(x_1,x_2)\,dx_1\,dx_2.
\end{equation}
For a concise introduction to the role of these weights,
see Vign\'eras \cite[\S2.2]{Vigneras1991}.
The exact definition of $v_M$ and the domains of validity are
part of Arthur's theory; \eqref{eq:arthur-weighted-pair}
should not be taken as a free-standing definition for arbitrary
singular $\gamma$. For $M=G$, no proper parabolic direction
remains, and the ordinary product-of-orbital-integrals pattern
reappears (with the local Weyl-discriminant normalization).

On the spectral side one obtains \emph{weighted characters}:
these are no longer just $\Theta_{\bar\pi}(f_1)\Theta_\pi(f_2)$.
They involve parabolically induced representations and
intertwining-operator data. In particular, taking the derivatives
of suitable operator families is a spectral analogue of keeping
track of the contributions at infinity. In general the resulting
weighted distributions for proper $M$ are not individually
conjugation-invariant; Arthur also develops invariant versions.

\subsection{The shape of Arthur's identity, and a comparison chart}
In a common schematic notation, Arthur's local formula is an
equality between sums of the form
\begin{equation}\label{eq:arthur-schematic}
 \underbrace{\sum_{M\in\mathcal L(M_0)} c_M
 \int_{\Gamma_{\mathrm{ell}}(M)}
 J_M(\gamma,f)\,d\gamma}_{\text{geometric expansion}}
 =
 \underbrace{\sum_{M\in\mathcal L(M_0)} c_M
 \int_{\mathscr T_{\mathrm{disc}}(M)}
 a^M(\tau)J_M(\tau,f)\,d\tau}_{\text{spectral expansion}},
\end{equation}
where the notation $\mathscr T_{\mathrm{disc}}(M)$ is a placeholder
for Arthur's discrete spectral parameter data, \emph{not} simply a
set of ordinary irreducible representations. The coefficients have
the usual Weyl-group/sign pattern
\[
 c_M=\frac{|W_0^M|}{|W_0^G|}
             (-1)^{\dim(\mathfrak a_M/\mathfrak a_G)},
\]
with measures and conventions as in Arthur's
\cite[Theorem 12.2]{ArthurLocal1991}.  This display deliberately
suppresses the detailed definitions of $J_M(\tau,f)$,
$a^M(\tau)$, parameter spaces, and normalizing factors. It is a
\emph{reading guide}, not an independent statement or proof of
Theorem 12.2. For $\mathrm{GL}_2$, the two Levi choices have
$c_G=1$ and $c_{M_0}=-\tfrac12$.

\begin{center}
\small
\begin{tabular}{@{}p{0.29\linewidth}p{0.59\linewidth}@{}}
\toprule
Finite-group step & Reductive $p$-adic replacement\\
\midrule
Counting measure on $G$ & Haar measure on a noncompact $G$\\
Every function is a test function & Hecke algebra $C_c^\infty(G)$\\
Finite sum over $\Irr(G)$ & Tempered Plancherel direct integral\\
Ordinary character $\chi_\pi(g)$ & Distribution character $\Theta_\pi(f)$\\
Conjugacy-class sum & Weyl integration over regular semisimple orbits\\
Kernel diagonal is summable & Kernel diagonal may diverge, even modulo $A_G$\\
No proper Levi terms & Levi-indexed geometric and spectral terms\\
Unweighted orbital products & Weighted orbital distributions $J_M(\gamma,f)$\\
Product of ordinary traces & Weighted characters $J_M(\tau,f)$\\
No truncation & Truncation and extraction of a regularized term\\
\bottomrule
\end{tabular}
\end{center}

\subsection{Suggested reading order and the limits of the analogy}
A practical route through the literature is to begin with
Vign\'eras \cite[\S1]{Vigneras1991} beside
Theorem~\ref{thm:diagonal} and Yang \cite{Yang2006}
beside Theorem~\ref{thm:global}. For the finite-geometric
variant, see Terras \cite[Chs.~22--23]{Terras1999}.
Then read \cite[\S\S1--2]{ArthurTowards1989} alongside
the diagonal-subgroup proof, and check the exact
$\mathrm{GL}_2$ divergence in Proposition~\ref{prop:gl2-divergence}.
Next read \cite[\S\S1--3]{ArthurLocal1991} for the kernel,
Plancherel expansion, and truncation; \cite[\S\S4--6]{ArthurLocal1991}
for the geometric weights; and \cite[\S\S7--12]{ArthurLocal1991}
for the spectral expansion and final identity.  The elementary
$p$-adic representation theory needed along the way is developed in
\cite{Casselman1995,Silberger1979}; \cite{Macdonald1971}
provides the spherical-function viewpoint.

\begin{remark}[What the finite model predicts, and what it does not]
The finite identities correctly predict a \emph{two-sided kernel},
a geometric/spectral comparison, and the importance of Weyl
integration and Plancherel theory.  They do \emph{not} predict the
analytic regularization, the family of parabolic directions, the
intertwining-operator terms, or the invariantization procedure.
Those are mathematical content, not bookkeeping corrections to the
finite formulas.
\end{remark}

\section{A brief literature review and reading path}
\label{sec:literature}
\paragraph{Four papers closest to the present approach.}
The finite-group identities in this note are classical; their significance
here is as a teaching device for the geometric/spectral dichotomy.
The most directly relevant references are the following.

\smallskip
\noindent\emph{Arthur's finite-group example (1989).}
In \emph{The Trace Formula and Hecke Operators}
\cite[\S1, Example~2, pp.~13--14]{ArthurHecke1989}, Arthur
specializes the compact-quotient Selberg formula to a finite group
$G$ and a subgroup $\Gamma$. Choosing a suitable irreducible
character as a test function shows that a spectral multiplicity
equals $\dim\pi^\Gamma$---the content of Frobenius reciprocity.
This is the closest direct precedent for our Theorem~\ref{thm:global}.
Arthur also explains why the compact trace formula can be viewed
as an extension of finite-group reciprocity to locally compact groups.

\smallskip
\noindent\emph{Vign\'eras's local analogue (1991).}
\emph{An Elementary Introduction to the Local Trace Formula of
J.~Arthur: The Case of Finite Groups} \cite{Vigneras1991}
begins with the $G\times G$-representation on $L[G]$ for a field $L$
of characteristic not dividing $|G|$. Her \S1 derives character
orthogonality, complete reducibility, and further structure of the
group algebra from a two-sided trace identity. The same
regular-bimodule mechanism underlies our
Theorem~\ref{thm:local} and the kernel computation in
Theorem~\ref{thm:diagonal}; our proof obtains it from the global
coset-space theorem by setting $(H,\Gamma)=(G\times G,\Delta G)$.
Vign\'eras then discusses elliptic elements, weighted orbital
integrals, and the actual reductive $p$-adic local formula in \S2,
with a final discussion of possible applications and limitations.
Her caution is worth emphasizing: a trace identity by itself
need not establish the character correspondences one would like.

\smallskip
\noindent\emph{Terras's finite Selberg analogues (1999).}
Chapter~22 of \emph{Fourier Analysis on Finite Groups and
Applications} \cite{Terras1999} develops finite versions of the
Selberg trace formula in relation to isospectral spectral geometry.
Chapter~23 studies the trace formula on finite upper half planes,
and Chapter~24 concerns a tree and Ihara's zeta function.
These chapters furnish geometric and spectral applications,
beyond the formal identity for $\C[\Gamma\backslash G]$.
The distinction matters: an elementary trace identity does not,
by itself, imply a graph-theoretic isospectrality theorem.

\smallskip
\noindent\emph{Yang's coset-space formulas (2006).}
Yang \cite{Yang2006} studies right-regular representations
on spaces of vector-valued functions associated with
$\Gamma\backslash G$ and states trace formulas building on
Vign\'eras. The scalar permutation module used here is a
simpler instance of the coset-space perspective.
Because Yang's paper appeared in 2006, it is a later related
contribution, not a source for the original 1996--2003 notes.

\paragraph{Finite representation theory.}
Serre's \emph{Linear Representations of Finite Groups} \cite{Serre}
supplies character orthogonality, induction, and Frobenius reciprocity.
In particular, it explains the averaging projection used in
\eqref{eq:mult} and the character-theoretic proof of
\eqref{eq:fourier}. For a trace-formula-first derivation
of these same foundations, compare Vign\'eras \cite[\S1]{Vigneras1991}.

\paragraph{The compact and noncompact global trace formulas.}
Arthur's later exposition \cite[Part~I, \S1]{ArthurIntro2005}
sets out the compact quotient and identifies the
centralizer-volume coefficient. Gelbart's lectures \cite{Gelbart1996}
introduce the general Arthur--Selberg framework, often through
$\mathrm{GL}_2$. Arthur's reductive-group papers
\cite{Arthur1978,Arthur1980,ArthurInvariant1981} develop the
truncation and invariant distributions needed once the quotient
is noncompact.

\paragraph{Background on reductive $p$-adic groups.}
Casselman \cite{Casselman1995} explains smooth and admissible
representations, parabolic induction, and intertwining operators;
Silberger \cite{Silberger1979} provides an analytic treatment.
Macdonald \cite{Macdonald1971} is useful for the Cartan
decomposition and spherical-function viewpoint used in the
$\mathrm{GL}_2(F)$ example of Section~\ref{sec:padic-bridge}.

\paragraph{The actual local trace formula.}
Arthur's \emph{Towards a Local Trace Formula}
\cite[\S\S1--2]{ArthurTowards1989} introduces the $G(F)\times G(F)$
kernel and discusses the compact case; Vign\'eras
\cite[\S2]{Vigneras1991} offers a particularly accessible
parallel introduction. Arthur's \emph{A Local Trace Formula}
\cite{ArthurLocal1991} develops truncation and the
Levi-indexed geometric and spectral expansions for the
noncompact reductive case. Those weighted terms are the
additional substance of the theory, not simply a notational
version of our finite sums.

\paragraph{Three guided reading routes.}
To understand \emph{the global finite analogy}, read
Theorem~\ref{thm:global}, Arthur's 1989 Example~2
\cite{ArthurHecke1989}, and Yang \cite{Yang2006}.
For \emph{the local finite analogy}, read
Theorem~\ref{thm:diagonal} beside Vign\'eras's \S1
\cite{Vigneras1991}. For \emph{geometric applications},
turn to Terras \cite[Chs.~22--24]{Terras1999}.
Only then proceed to Section~\ref{sec:padic-bridge},
Vign\'eras's \S2, and Arthur's local papers.

\section{Glossary of terms and notation}
\begin{description}[style=nextline,leftmargin=1.4em,itemsep=4pt]
\item[$A_M$ and $\mathfrak a_M$]
The maximal $F$-split torus in the center of a Levi subgroup
$M$ and its associated real vector space of logarithmic
rational-character data, respectively.
\item[Arthur--Selberg trace formula]
A family of identities equating geometric orbital terms with spectral
representation-theoretic terms, often after truncation.
\item[Character $\chi_\pi$]
The class function $g\mapsto\tr(\pi(g))$ of a representation $\pi$.
\item[Class function]
A function constant on conjugacy classes.
\item[Compact quotient]
A homogeneous space $\Gamma\backslash H$ which is compact, even if
$H$ itself is not compact.
\item[Conjugacy class $\operatorname{Cl}_G(\gamma)$]
The set $\{x^{-1}\gamma x:x\in G\}$.
\item[Centralizer $G_\gamma$]
The subgroup of elements commuting with $\gamma$; if
$\gamma\in\Gamma$ then $\Gamma_\gamma=G_\gamma\cap\Gamma$.
\item[Contragredient representation $\bar\pi$]
The dual representation; for a finite group one may choose a unitary
model in which $\chi_{\bar\pi}(g)=\overline{\chi_\pi(g)}$.
\item[Counting measure]
The convention $\int_X f=\sum_{x\in X}f(x)$ for finite $X$.
\item[Cartan double cosets]
In the $\mathrm{GL}_2(F)$ example, the sets
$\bar K\bar a_n\bar K$ partition $Z\backslash G$;
their measure grows as $(q+1)q^{n-1}$.
\item[Diagonal subgroup $\Delta G$]
The subgroup $\{(g,g):g\in G\}$ of $G\times G$.
\item[Eisenstein series]
Functions built from parabolic induction that help describe continuous
automorphic spectrum on noncompact arithmetic quotients.
\item[Frobenius reciprocity]
The adjunction relating $\Hom_G(\Ind_\Gamma^G\sigma,\pi)$ to
$\Hom_\Gamma(\sigma,\Res_\Gamma^G\pi)$.
\item[Geometric side]
The trace expansion organized by conjugacy classes and orbital integrals.
\item[Harish-Chandra's Plancherel theorem]
Describes the regular representation of a reductive local group
by a direct integral of tempered representations and a
Plancherel measure.
\item[Induced representation]
A representation of $G$ built from one of a subgroup, written
$\Ind_\Gamma^G\sigma$.
\item[Invariant vectors $\pi^\Gamma$]
Vectors fixed by every element of $\Gamma$; their dimension is the
multiplicity in \eqref{eq:mult}.
\item[Irreducible representation]
A representation with no nonzero proper invariant subspace.
\item[Hecke algebra $\mathcal H(G)$]
The convolution algebra of compactly supported locally constant
functions on a totally disconnected group $G$.
\item[Levi subgroup]
A reductive factor of a parabolic subgroup; Levi subgroups index many
terms in Arthur's noncompact global and local expansions.
\item[Local field]
A locally compact nondiscrete topological field, such as
$\mathbb R$, $\C$, or a finite extension of $\mathbb Q_p$.
\item[Normalized parabolic induction]
Induction from a Levi subgroup of a parabolic, including the factor
$\delta_P^{1/2}$ built from the modulus character.
\item[Orbital integral / orbital sum]
An integral (or finite sum) of a test function along a conjugacy orbit.
Here $\Orb_G$ is deliberately the \emph{sum over all $x\in G$};
it includes a factor of $|G_\gamma|$ relative to summing once
per element of $G_\gamma\backslash G$.
\item[Peter--Weyl decomposition]
A decomposition of the regular representation of a compact group into
matrix coefficients of irreducibles.  For finite $G$ it yields
\eqref{eq:peter-weyl}.
\item[Schur orthogonality]
The orthogonality relations for irreducible matrix coefficients and,
in particular, for irreducible characters.
\item[Smooth admissible representation]
A smooth representation for which the fixed-vector space under
any compact open subgroup is finite-dimensional.
\item[Spectral side]
The trace expansion organized by irreducible representations and
multiplicities or spectral measure.
\item[Tempered representation]
A representation in the Plancherel support, broadly analogous to
the part of Fourier analysis that contributes to the regular
$L^2$-spectrum; precise definitions use matrix-coefficient decay.
\item[Test function]
A function $f$ used to form a convolution operator; in finite groups
it is any element of $\C[G]$.
\item[Truncation]
A regularization procedure used when the kernel or its diagonal is
not integrable on a noncompact space.
\item[Unipotent radical $N$]
The unipotent normal subgroup in a parabolic decomposition $P=MN$.
\item[Weighted orbital integral / weighted character]
A more elaborate geometric or spectral distribution appearing in
Arthur's local/global trace formulas for noncompact reductive groups.
\item[Weyl integration formula]
A decomposition of integrals into orbital integrals. In the finite
setting it is \eqref{eq:weyl}.
\end{description}

\section{Exercises and optional extensions}
\begin{exercise}[A general $S_3$ calculation]
For $a,b,c\in\C$, set
$f=a f_1+b f_2+c f_3$, using the normalized class indicators
in \eqref{eq:basis-orb} for $S_3$. Verify directly that
\[
 \sum_{\pi\in\Irr(S_3)}|\chi_\pi(f)|^2
 =\frac{|a|^2}{6}+\frac{|b|^2}{2}+\frac{|c|^2}{3}.
\]
\end{exercise}
\begin{exercise}[A nonnormal subgroup]
For $G=S_3$ and $\Gamma=\langle(12)\rangle$,
check the trace formula on the three point masses
$\delta_e$, $\delta_{(12)}$, and $\delta_{(123)}$.
In the last case, the trace is zero.
\end{exercise}
\begin{exercise}[A diagonal-global computation]
Apply Theorem~\ref{thm:global} directly to
$(G\times G,\Delta G)$ with $f_1=\delta_u$ and
$f_2=\delta_v$. Recover column orthogonality
\eqref{eq:columnorth}.
\end{exercise}

\begin{exercise}[The $\mathrm{GL}_2$ shell count]
For a finite ring $R=\mathcal O_F/\varpi^n\mathcal O_F$,
verify that $\#\mathbb P^1(R)=q^n+q^{n-1}$. Deduce
$[K:K\cap a_nKa_n^{-1}]=(q+1)q^{n-1}$.
\end{exercise}
\begin{exercise}[Two kinds of elliptic elements]
In $\mathrm{GL}_2(F)$, compare the centralizers of
$\operatorname{diag}(a,d)$, $a\ne d$, and a regular matrix with
irreducible quadratic characteristic polynomial. Explain why
one class is elliptic only in the diagonal Levi and the other is
elliptic modulo the center in $G$.
\end{exercise}

\begin{exercise}[Arthur's Frobenius-reciprocity example]
Fix $\pi\in\Irr(G)$ and set
$f_\pi(g)=\chi_\pi(g^{-1})$. Use Theorem~\ref{thm:global}
and the orthogonality of irreducible characters to isolate
$m_\Gamma(\pi)$. Recover
$m_\Gamma(\pi)=\dim\pi^\Gamma$ and compare the calculation
with \cite[Example~2]{ArthurHecke1989}.
\end{exercise}
\begin{exercise}[Reading Vign\'eras and Terras]
Compare Vign\'eras's characteristic-sum notation
\cite[\S1.1]{Vigneras1991} with
$\Orb_G(f,\gamma)$ in \eqref{eq:def-orbit}.
Then consult Terras \cite[Ch.~22]{Terras1999}:
what finite geometric or spectral data enter beyond the
permutation module $\C[\Gamma\backslash G]$?
\end{exercise}

\appendix
\section{A short exact sequence and group cohomology}
The conjugation action of $G$ on $\cG$ is
$(h\cdot f)(x)=f(h^{-1}xh)$. Its averaging projection is
\[
 (Pf)(x)=\frac1{|G|}\sum_{h\in G}f(h^{-1}xh)
        =\frac1{|G|}\Orb_G(f,x).
\]
It satisfies $P^2=P$ and has image $\cG^G$, the class functions.
Moreover,
\[
 \ker P=\operatorname{span}_\C
 \{f-h\cdot f:f\in\cG,\ h\in G\}=:I(G).
\]
Consequently
\[
 0\longrightarrow I(G)\longrightarrow\cG
   \overset P\longrightarrow\cG^G\longrightarrow0
\]
is a split exact sequence of complex $G$-modules.  The word
``span'' is necessary: the \emph{set} of all single differences
$f-h\cdot f$ need not itself be a linear subspace.
Since $G$ is finite and the coefficient field has characteristic
zero, averaging makes the invariants functor exact. In particular,
for every complex $G$-module $M$ and every $n>0$,
\[
 H^n(G,M)=0.
\]
Thus all higher cohomology groups of $I(G)$, $\cG$, and
$\cG^G$ vanish. The nontrivial information here is in degree zero:
$\cG^G$ is the class-function space and $I(G)^G=0$.

\section{Summary of substantive corrections to the original notes}
\begin{enumerate}[leftmargin=*,itemsep=5pt]
\item Corrected $S_3$'s standard representation from the wrong
quotient $\C^3/\{z_1+z_2+z_3=0\}$ to the two-dimensional
subspace $\{z_1+z_2+z_3=0\}$ (or quotient by the diagonal line).
\item Distinguished point masses from class indicators in the
character-table argument; the latter have a class-size factor.
\item Distinguished the normalized character inner product from
unnormalized point masses, and noted that characters span class
functions rather than all functions on $G$.
\item Replaced the inverse-argument right translation, which was an
anti-representation under ordinary composition, by the correct
representation $R(g)\varphi(x)=\varphi(xg)$.
\item Recomputed the kernel with counting measure on
$\Gamma\backslash G$; fixed the factors $|\Gamma|$ and $|G|$.
\item Corrected the geometric side: sum over $\Gamma$-conjugacy
classes with $1/|\Gamma_\gamma|$, or sum over $G$-classes with
$|[\gamma]_G\cap\Gamma|/|\Gamma|$.
\item Replaced the faulty extraction of spectral multiplicities
by the projection formula $m_\Gamma(\pi)=\dim\pi^\Gamma$.
\item Supplied the missing $1/|\Gamma|$ in the induced-character
formula and removed the unjustified assumption that every subgroup
is stable under conjugation.
\item Replaced the index-two argument by the direct identity
$\Ind_\Gamma^G\id=\id\oplus\eps$.
\item Corrected the $A_5$ notation: $(1\pm\sqrt5)/2$ are real
algebraic conjugates, not complex conjugates.
\item Corrected the character-table matrix identity to
$A^*A=\diag(|G_{\gamma_j}|^{-1})$, not $AA^*$.
\item Kept the valid bilinear local identity, strengthened it to
complex functions with a positive Hermitian form, and identified it
as the global trace formula for $(G\times G,\Delta G)$.
\item Distinguished the compact local analogue from the full
Arthur local trace formula, which requires truncation and weighted
terms in noncompact situations.
\item In the subsequent expansion, showed explicitly for
$\mathrm{GL}_2(F)$ that the untruncated local diagonal integral
diverges even modulo the center and detailed the role of Levi groups,
parabolic induction, and weighted distributions.
\end{enumerate}

\section{Concluding perspective}
The finite trace formulas rest on two familiar operations:
computing the trace of a convolution operator on its diagonal, and
computing the trace after decomposing a representation into
irreducibles.  For $\C[\Gamma\backslash G]$ this gives the
compact-global identity; for the two-sided action of $G\times G$ on
$\C[G]$ it gives the compact-local identity.  Their common structure
is especially transparent because the second is literally the first
applied to the diagonal subgroup.  The main conceptual leap in
Selberg's and Arthur's work is to make related comparisons meaningful
for noncompact groups and quotients, where continuous spectrum,
parabolic induction, and truncation all become central. The precise
$\mathrm{GL}_2(F)$ calculation in Section~\ref{sec:padic-bridge}
shows why the passage to the local trace formula is not merely
a change from finite sums to integrals.

\begin{thebibliography}{99}
\bibitem{ArthurHecke1989}
J. Arthur, ``The trace formula and Hecke operators,'' in
\emph{Number Theory, Trace Formulas and Discrete Groups},
Academic Press,
1989, pp.~11--27.
\url{https://www.claymath.org/library/cw/arthur/pdf/31.pdf}

\bibitem{Vigneras1991}
M.-F. Vign\'eras, ``An elementary introduction to the local
trace formula of J.~Arthur: the case of finite groups,''
in a Deutsche Mathematiker-Vereinigung jubilee volume,
B.~G. Teubner, Stuttgart, 1991.
\url{https://perso.imj-prg.fr/mariefrance-vigneras/wp-content/uploads/vigneras-pub/dmv.pdf}

\bibitem{Terras1999}
A. Terras, \emph{Fourier Analysis on Finite Groups and
Applications}, London Mathematical Society Student Texts 43,
Cambridge University Press, 1999; esp. Ch.~22,
``Selberg's trace formula and isospectral non-isomorphic
graphs,'' pp.~385--393, and Ch.~23, pp.~394--407.
\url{https://doi.org/10.1017/CBO9780511626265.025}

\bibitem{Yang2006}
J.-H. Yang, ``Trace formulas on finite groups,''
\emph{Communications of the Korean Mathematical Society}
\textbf{21} (2006), no.~1, 17--25.
\url{https://arxiv.org/abs/math/0612362}

\bibitem{Casselman1995}
W. Casselman, \emph{Introduction to the Theory of Admissible
Representations of $p$-Adic Reductive Groups},
draft, 1 May 1995.
\url{https://www.math.utah.edu/~ptrapa/math-library/casselman/casselman-p-adic-book.pdf}

\bibitem{Silberger1979}
A. J. Silberger, \emph{Introduction to Harmonic Analysis on
Reductive $p$-Adic Groups}, Mathematical Notes 23, Princeton
University Press, 1979.

\bibitem{Macdonald1971}
I. G. Macdonald, \emph{Spherical Functions on a Group of
$p$-Adic Type}, Publications of the Ramanujan Institute,
University of Madras, no. 2, 1971.
\url{https://math.soimeme.org/~arunram/Resources/SFOAGOPTContents.html}

\bibitem{Serre}
J.-P. Serre, \emph{Linear Representations of Finite Groups},
Graduate Texts in Mathematics 42, Springer, New York, 1977.

\bibitem{Arthur1978}
J. Arthur, ``A trace formula for reductive groups I:
Terms associated to classes in $G(\mathbb Q)$,''
\emph{Duke Math. J.} \textbf{45} (1978), 911--952.

\bibitem{Arthur1980}
J. Arthur, ``A trace formula for reductive groups II:
Applications of a truncation operator,''
\emph{Compositio Math.} \textbf{40} (1980), 87--121.

\bibitem{ArthurInvariant1981}
J. Arthur, ``The trace formula in invariant form,''
\emph{Ann. of Math.} (2) \textbf{114} (1981), 1--74.

\bibitem{ArthurTowards1989}
J. Arthur, ``Towards a local trace formula,'' in
\emph{Algebraic Analysis, Geometry, and Number Theory},
Johns Hopkins University Press, Baltimore, 1989, pp. 1--23.
\url{https://www.claymath.org/library/cw/arthur/pdf/36.pdf}

\bibitem{ArthurLocal1991}
J. Arthur, ``A local trace formula,''
\emph{Publications Math\'ematiques de l'IH\'ES} \textbf{73}
(1991), 5--96. \url{https://doi.org/10.1007/BF02699256}

\bibitem{Gelbart1996}
S. Gelbart, \emph{Lectures on the Arthur--Selberg Trace Formula},
University Lecture Series 9, American Mathematical Society,
Providence, RI, 1996. \url{https://arxiv.org/abs/math/9505206}

\bibitem{ArthurIntro2005}
J. Arthur, ``An introduction to the trace formula,'' in
\emph{Harmonic Analysis, the Trace Formula, and Shimura Varieties},
Clay Mathematics Proceedings 4, American Mathematical Society,
Providence, RI, 2005, pp. 1--263.
\url{https://www.claymath.org/online-resources/arthur/}
\end{thebibliography}
\end{document}
